% The package's Python surface, typeset (issue #576). The entries themselves
% are generated by infra/api_map.py into generated/api_map.tex, which this
% document inputs; nothing here is written by hand except the macros the
% generator emits against and the framing below.
\documentclass[11pt,a4paper]{article}
% Packages every document in this directory shares. Kept apart from
% notation.tex so a symbol change and a package change are separate reviews.
\usepackage[margin=1in]{geometry}
\usepackage{amsmath}
\usepackage{amssymb}
\usepackage{booktabs}
\usepackage{graphicx}
\usepackage{tikz}
\usetikzlibrary{shapes.geometric}
\usepackage{microtype}
\usepackage{algorithm}
\usepackage{algpseudocode}
\usepackage[numbers]{natbib}
\usepackage[hidelinks]{hyperref}

% The one visual vocabulary the problem-statement sketches share. Circles are
% variables, squares are factors, shading marks where an observation enters,
% and a triangle marks a factor that gates on a second variable. Each of the
% eleven sketches carried a private copy, between them four circle diameters,
% three square sizes and four fonts, so a reader learned the convention once
% per figure (issue #529). A figure that needs a shape of its own adds it
% locally, on top of these.
\tikzset{
  fgvariable/.style={circle, draw, fill=white, minimum size=0.58cm,
                     inner sep=0.5pt, font=\footnotesize},
  fgobserved/.style={fgvariable, fill=black!12},
  fgfactor/.style={rectangle, draw, fill=white, minimum size=0.16cm,
                   inner sep=0pt},
  fgdata/.style={fgfactor, fill=black!25},
  fgindicator/.style={fgfactor, fill=black!55},
  fggate/.style={regular polygon, regular polygon sides=3, draw, fill=white,
                 minimum size=0.25cm, inner sep=0pt},
  fglink/.style={draw, thick},
  fgabsent/.style={draw, thin, gray, dashed},
  fgnote/.style={anchor=west, font=\footnotesize}
}

% Shared notation for every document in this directory.
%
% One definition per symbol, input by both the paper and the textbook, so a
% symbol cannot mean two things in two places (issue #249). A document that
% needs a symbol adds it here rather than defining its own; the notation table
% in the textbook is generated from this list by hand and reviewed against it.

% --- Structures ----------------------------------------------------------
\newcommand{\tree}{\mathcal{T}}          % a tree topology
\newcommand{\graph}{\mathcal{G}}         % a general undirected graph
\newcommand{\states}{\mathcal{A}}        % the state alphabet
\newcommand{\rootnode}{\rho}             % the root of a rooted tree
\newcommand{\neighbourhood}{\mathcal{N}} % a discrete move neighbourhood

% --- Continuous parameters ------------------------------------------------
\newcommand{\Q}{\mathbf{Q}}              % a rate matrix
\newcommand{\Pmat}{\mathbf{P}}           % a transition probability matrix
\newcommand{\SubMat}{\mathbf{S}}         % a symmetric exchangeability matrix
\newcommand{\bpi}{\boldsymbol{\pi}}      % a stationary or initial distribution
\newcommand{\branch}{\boldsymbol{\tau}}  % branch lengths
\newcommand{\params}{\boldsymbol{\theta}}% unconstrained parameters
\newcommand{\constrained}{\phi}          % the constraint map

% --- Data and objectives --------------------------------------------------
\newcommand{\data}{\mathbf{D}}           % observed data
\newcommand{\objective}{\mathcal{L}}     % the scalar being optimized
\newcommand{\energy}{H}                  % a Hamiltonian
\newcommand{\hidden}{\mathbf{Z}}         % a hidden state path
\newcommand{\emission}{\mathcal{E}}      % an emission family

% --- Reading a generated caption -----------------------------------------
% The captions under figures/ are written by the QA build, never by hand
% (docs/CLAUDE.md). This reads one in so a document quotes it rather than
% restating it.
\newread\qacaptionfile
\newcommand{\qacaptionread}[2]{%
  \begingroup
  \endlinechar=-1
  \openin\qacaptionfile=#1\relax
  \global\read\qacaptionfile to #2
  \closein\qacaptionfile
  \endgroup
}


% One row per entry: the name in the left column and its own summary line in
% the right. A `longtable` rather than a `tabular` because a module's surface
% runs over a page break, and the surrounding text is what decides where a
% module starts.
\newenvironment{apitable}
  {\begin{longtable}{@{}p{0.34\textwidth}p{0.62\textwidth}@{}}}
  {\end{longtable}}
\newcommand{\apiclass}[2]{\texttt{#1} & \emph{class.} #2 \\}
\newcommand{\apifunction}[2]{\texttt{#1} & #2 \\}
\newcommand{\apimethod}[2]{\quad\texttt{#1} & #2 \\}

\title{The Package's Python Surface \\
\large Every module, class, function and public method, with the summary line
it carries in its own docstring}
\author{M.\ J.\ Wilson}
\date{\today}

\begin{document}
\maketitle

% The map as a whole. Every label the generated entries carry hangs under it
% --- `sec:api:reading`, `sec:api:package:<package>`, `sec:api:<module>` ---
% and code citing one of those names this stem, the way `app:algorithms` in
% the textbook is the stem of `app:algorithms:samplers`. Written here and not
% into `generated/api_map.tex`, which is not committed and so is not in the
% tree a citation is read against. `\phantomsection` because the point being
% named is the document and not a numbered sectional unit.
\phantomsection
\label{sec:api}

\begin{abstract}
\noindent
A map of what the package exposes, generated from the sources rather than
written: `infra/api\_map.py` reads each module by `ast` --- never by importing
it, so no import side effect and no compiled extension can drop a module from
the map --- and lays out the first paragraph of every docstring it finds. The
companion documents state the mathematics (the textbook) and the measurements
(the paper); this one states the surface, and it is the only one of the three
whose content no one edits.
\end{abstract}

\tableofcontents
\clearpage

% The package's Python surface, typeset (issue #576). The entries themselves
% are generated by infra/api_map.py into generated/api_map.tex, which this
% document inputs; nothing here is written by hand except the macros the
% generator emits against and the framing below.
\documentclass[11pt,a4paper]{article}
% Packages every document in this directory shares. Kept apart from
% notation.tex so a symbol change and a package change are separate reviews.
\usepackage[margin=1in]{geometry}
\usepackage{amsmath}
\usepackage{amssymb}
\usepackage{booktabs}
\usepackage{graphicx}
\usepackage{tikz}
\usetikzlibrary{shapes.geometric}
\usepackage{microtype}
\usepackage{algorithm}
\usepackage{algpseudocode}
\usepackage[numbers]{natbib}
\usepackage[hidelinks]{hyperref}

% The one visual vocabulary the problem-statement sketches share. Circles are
% variables, squares are factors, shading marks where an observation enters,
% and a triangle marks a factor that gates on a second variable. Each of the
% eleven sketches carried a private copy, between them four circle diameters,
% three square sizes and four fonts, so a reader learned the convention once
% per figure (issue #529). A figure that needs a shape of its own adds it
% locally, on top of these.
\tikzset{
  fgvariable/.style={circle, draw, fill=white, minimum size=0.58cm,
                     inner sep=0.5pt, font=\footnotesize},
  fgobserved/.style={fgvariable, fill=black!12},
  fgfactor/.style={rectangle, draw, fill=white, minimum size=0.16cm,
                   inner sep=0pt},
  fgdata/.style={fgfactor, fill=black!25},
  fgindicator/.style={fgfactor, fill=black!55},
  fggate/.style={regular polygon, regular polygon sides=3, draw, fill=white,
                 minimum size=0.25cm, inner sep=0pt},
  fglink/.style={draw, thick},
  fgabsent/.style={draw, thin, gray, dashed},
  fgnote/.style={anchor=west, font=\footnotesize}
}

% Shared notation for every document in this directory.
%
% One definition per symbol, input by both the paper and the textbook, so a
% symbol cannot mean two things in two places (issue #249). A document that
% needs a symbol adds it here rather than defining its own; the notation table
% in the textbook is generated from this list by hand and reviewed against it.

% --- Structures ----------------------------------------------------------
\newcommand{\tree}{\mathcal{T}}          % a tree topology
\newcommand{\graph}{\mathcal{G}}         % a general undirected graph
\newcommand{\states}{\mathcal{A}}        % the state alphabet
\newcommand{\rootnode}{\rho}             % the root of a rooted tree
\newcommand{\neighbourhood}{\mathcal{N}} % a discrete move neighbourhood

% --- Continuous parameters ------------------------------------------------
\newcommand{\Q}{\mathbf{Q}}              % a rate matrix
\newcommand{\Pmat}{\mathbf{P}}           % a transition probability matrix
\newcommand{\SubMat}{\mathbf{S}}         % a symmetric exchangeability matrix
\newcommand{\bpi}{\boldsymbol{\pi}}      % a stationary or initial distribution
\newcommand{\branch}{\boldsymbol{\tau}}  % branch lengths
\newcommand{\params}{\boldsymbol{\theta}}% unconstrained parameters
\newcommand{\constrained}{\phi}          % the constraint map

% --- Data and objectives --------------------------------------------------
\newcommand{\data}{\mathbf{D}}           % observed data
\newcommand{\objective}{\mathcal{L}}     % the scalar being optimized
\newcommand{\energy}{H}                  % a Hamiltonian
\newcommand{\hidden}{\mathbf{Z}}         % a hidden state path
\newcommand{\emission}{\mathcal{E}}      % an emission family

% --- Reading a generated caption -----------------------------------------
% The captions under figures/ are written by the QA build, never by hand
% (docs/CLAUDE.md). This reads one in so a document quotes it rather than
% restating it.
\newread\qacaptionfile
\newcommand{\qacaptionread}[2]{%
  \begingroup
  \endlinechar=-1
  \openin\qacaptionfile=#1\relax
  \global\read\qacaptionfile to #2
  \closein\qacaptionfile
  \endgroup
}


% One row per entry: the name in the left column and its own summary line in
% the right. A `longtable` rather than a `tabular` because a module's surface
% runs over a page break, and the surrounding text is what decides where a
% module starts.
\newenvironment{apitable}
  {\begin{longtable}{@{}p{0.34\textwidth}p{0.62\textwidth}@{}}}
  {\end{longtable}}
\newcommand{\apiclass}[2]{\texttt{#1} & \emph{class.} #2 \\}
\newcommand{\apifunction}[2]{\texttt{#1} & #2 \\}
\newcommand{\apimethod}[2]{\quad\texttt{#1} & #2 \\}

\title{The Package's Python Surface \\
\large Every module, class, function and public method, with the summary line
it carries in its own docstring}
\author{M.\ J.\ Wilson}
\date{\today}

\begin{document}
\maketitle

% The map as a whole. Every label the generated entries carry hangs under it
% --- `sec:api:reading`, `sec:api:package:<package>`, `sec:api:<module>` ---
% and code citing one of those names this stem, the way `app:algorithms` in
% the textbook is the stem of `app:algorithms:samplers`. Written here and not
% into `generated/api_map.tex`, which is not committed and so is not in the
% tree a citation is read against. `\phantomsection` because the point being
% named is the document and not a numbered sectional unit.
\phantomsection
\label{sec:api}

\begin{abstract}
\noindent
A map of what the package exposes, generated from the sources rather than
written: `infra/api\_map.py` reads each module by `ast` --- never by importing
it, so no import side effect and no compiled extension can drop a module from
the map --- and lays out the first paragraph of every docstring it finds. The
companion documents state the mathematics (the textbook) and the measurements
(the paper); this one states the surface, and it is the only one of the three
whose content no one edits.
\end{abstract}

\tableofcontents
\clearpage

% The package's Python surface, typeset (issue #576). The entries themselves
% are generated by infra/api_map.py into generated/api_map.tex, which this
% document inputs; nothing here is written by hand except the macros the
% generator emits against and the framing below.
\documentclass[11pt,a4paper]{article}
% Packages every document in this directory shares. Kept apart from
% notation.tex so a symbol change and a package change are separate reviews.
\usepackage[margin=1in]{geometry}
\usepackage{amsmath}
\usepackage{amssymb}
\usepackage{booktabs}
\usepackage{graphicx}
\usepackage{tikz}
\usetikzlibrary{shapes.geometric}
\usepackage{microtype}
\usepackage{algorithm}
\usepackage{algpseudocode}
\usepackage[numbers]{natbib}
\usepackage[hidelinks]{hyperref}

% The one visual vocabulary the problem-statement sketches share. Circles are
% variables, squares are factors, shading marks where an observation enters,
% and a triangle marks a factor that gates on a second variable. Each of the
% eleven sketches carried a private copy, between them four circle diameters,
% three square sizes and four fonts, so a reader learned the convention once
% per figure (issue #529). A figure that needs a shape of its own adds it
% locally, on top of these.
\tikzset{
  fgvariable/.style={circle, draw, fill=white, minimum size=0.58cm,
                     inner sep=0.5pt, font=\footnotesize},
  fgobserved/.style={fgvariable, fill=black!12},
  fgfactor/.style={rectangle, draw, fill=white, minimum size=0.16cm,
                   inner sep=0pt},
  fgdata/.style={fgfactor, fill=black!25},
  fgindicator/.style={fgfactor, fill=black!55},
  fggate/.style={regular polygon, regular polygon sides=3, draw, fill=white,
                 minimum size=0.25cm, inner sep=0pt},
  fglink/.style={draw, thick},
  fgabsent/.style={draw, thin, gray, dashed},
  fgnote/.style={anchor=west, font=\footnotesize}
}

% Shared notation for every document in this directory.
%
% One definition per symbol, input by both the paper and the textbook, so a
% symbol cannot mean two things in two places (issue #249). A document that
% needs a symbol adds it here rather than defining its own; the notation table
% in the textbook is generated from this list by hand and reviewed against it.

% --- Structures ----------------------------------------------------------
\newcommand{\tree}{\mathcal{T}}          % a tree topology
\newcommand{\graph}{\mathcal{G}}         % a general undirected graph
\newcommand{\states}{\mathcal{A}}        % the state alphabet
\newcommand{\rootnode}{\rho}             % the root of a rooted tree
\newcommand{\neighbourhood}{\mathcal{N}} % a discrete move neighbourhood

% --- Continuous parameters ------------------------------------------------
\newcommand{\Q}{\mathbf{Q}}              % a rate matrix
\newcommand{\Pmat}{\mathbf{P}}           % a transition probability matrix
\newcommand{\SubMat}{\mathbf{S}}         % a symmetric exchangeability matrix
\newcommand{\bpi}{\boldsymbol{\pi}}      % a stationary or initial distribution
\newcommand{\branch}{\boldsymbol{\tau}}  % branch lengths
\newcommand{\params}{\boldsymbol{\theta}}% unconstrained parameters
\newcommand{\constrained}{\phi}          % the constraint map

% --- Data and objectives --------------------------------------------------
\newcommand{\data}{\mathbf{D}}           % observed data
\newcommand{\objective}{\mathcal{L}}     % the scalar being optimized
\newcommand{\energy}{H}                  % a Hamiltonian
\newcommand{\hidden}{\mathbf{Z}}         % a hidden state path
\newcommand{\emission}{\mathcal{E}}      % an emission family

% --- Reading a generated caption -----------------------------------------
% The captions under figures/ are written by the QA build, never by hand
% (docs/CLAUDE.md). This reads one in so a document quotes it rather than
% restating it.
\newread\qacaptionfile
\newcommand{\qacaptionread}[2]{%
  \begingroup
  \endlinechar=-1
  \openin\qacaptionfile=#1\relax
  \global\read\qacaptionfile to #2
  \closein\qacaptionfile
  \endgroup
}


% One row per entry: the name in the left column and its own summary line in
% the right. A `longtable` rather than a `tabular` because a module's surface
% runs over a page break, and the surrounding text is what decides where a
% module starts.
\newenvironment{apitable}
  {\begin{longtable}{@{}p{0.34\textwidth}p{0.62\textwidth}@{}}}
  {\end{longtable}}
\newcommand{\apiclass}[2]{\texttt{#1} & \emph{class.} #2 \\}
\newcommand{\apifunction}[2]{\texttt{#1} & #2 \\}
\newcommand{\apimethod}[2]{\quad\texttt{#1} & #2 \\}

\title{The Package's Python Surface \\
\large Every module, class, function and public method, with the summary line
it carries in its own docstring}
\author{M.\ J.\ Wilson}
\date{\today}

\begin{document}
\maketitle

% The map as a whole. Every label the generated entries carry hangs under it
% --- `sec:api:reading`, `sec:api:package:<package>`, `sec:api:<module>` ---
% and code citing one of those names this stem, the way `app:algorithms` in
% the textbook is the stem of `app:algorithms:samplers`. Written here and not
% into `generated/api_map.tex`, which is not committed and so is not in the
% tree a citation is read against. `\phantomsection` because the point being
% named is the document and not a numbered sectional unit.
\phantomsection
\label{sec:api}

\begin{abstract}
\noindent
A map of what the package exposes, generated from the sources rather than
written: `infra/api\_map.py` reads each module by `ast` --- never by importing
it, so no import side effect and no compiled extension can drop a module from
the map --- and lays out the first paragraph of every docstring it finds. The
companion documents state the mathematics (the textbook) and the measurements
(the paper); this one states the surface, and it is the only one of the three
whose content no one edits.
\end{abstract}

\tableofcontents
\clearpage

% The package's Python surface, typeset (issue #576). The entries themselves
% are generated by infra/api_map.py into generated/api_map.tex, which this
% document inputs; nothing here is written by hand except the macros the
% generator emits against and the framing below.
\documentclass[11pt,a4paper]{article}
\input{preamble}
\input{notation}

% One row per entry: the name in the left column and its own summary line in
% the right. A `longtable` rather than a `tabular` because a module's surface
% runs over a page break, and the surrounding text is what decides where a
% module starts.
\newenvironment{apitable}
  {\begin{longtable}{@{}p{0.34\textwidth}p{0.62\textwidth}@{}}}
  {\end{longtable}}
\newcommand{\apiclass}[2]{\texttt{#1} & \emph{class.} #2 \\}
\newcommand{\apifunction}[2]{\texttt{#1} & #2 \\}
\newcommand{\apimethod}[2]{\quad\texttt{#1} & #2 \\}

\title{The Package's Python Surface \\
\large Every module, class, function and public method, with the summary line
it carries in its own docstring}
\author{M.\ J.\ Wilson}
\date{\today}

\begin{document}
\maketitle

% The map as a whole. Every label the generated entries carry hangs under it
% --- `sec:api:reading`, `sec:api:package:<package>`, `sec:api:<module>` ---
% and code citing one of those names this stem, the way `app:algorithms` in
% the textbook is the stem of `app:algorithms:samplers`. Written here and not
% into `generated/api_map.tex`, which is not committed and so is not in the
% tree a citation is read against. `\phantomsection` because the point being
% named is the document and not a numbered sectional unit.
\phantomsection
\label{sec:api}

\begin{abstract}
\noindent
A map of what the package exposes, generated from the sources rather than
written: `infra/api\_map.py` reads each module by `ast` --- never by importing
it, so no import side effect and no compiled extension can drop a module from
the map --- and lays out the first paragraph of every docstring it finds. The
companion documents state the mathematics (the textbook) and the measurements
(the paper); this one states the surface, and it is the only one of the three
whose content no one edits.
\end{abstract}

\tableofcontents
\clearpage

\input{generated/api_map}

\end{document}


\end{document}


\end{document}


\end{document}
